\fichatitulo{30 · REVISÃO}{Preprint v1 · 2023}{A Survey of Large Language Models}
{\small Zhao et al. \textit{A Survey of Large Language Models}. Preprint v1 · 2023. \href{https://arxiv.org/pdf/2303.18223v1}{Fonte consultada}.\par}

\campo{Problema e objetivo}
Quais etapas organizam o desenvolvimento e uso dos LLMs?

\campo{Principal contribuição · síntese interpretativa}
Articula pré-treinamento, adaptação, utilização e avaliação de capacidades em um panorama técnico.

\campo{Método / natureza da evidência}
Revisão de modelos, estratégias e tarefas, com discussão de desafios e direções de pesquisa.

\campo{Resultado ou argumento principal}
Permite distinguir formação do modelo, adaptação posterior e uso por prompts; também discute sensibilidade às instruções.

\campo{Excertos-chave · seleção literal}
\begin{itemize}
\excerto{1}{pre-training, adaptation tuning, utilization, and capacity evaluation}{Resumo.}
\excerto{2}{prompt sensitivity}{Discussão de avaliação.}
\end{itemize}

\campo{Limitações e análise crítica}
Análise da ficha: conceitos e exemplos refletem a versão de março de 2023. Descrições de capacidades emergentes não devem ser convertidas em garantias universais de comportamento.

\campo{Aproveitamento na dissertação · proposta}
Fundamentação: delimitar quais etapas são modificadas pela dissertação e quais permanecem herdadas do modelo.

\registro{Fonte consultada nesta preparação: Resumo, introdução e trechos de alinhamento, avaliação e perspectivas. Chave: \texttt{zhaoEtAl2023llmSurvey}.}
